\documentclass{article}
\usepackage{amsmath}
\usepackage{graphicx}
\graphicspath{ {images/} }
\usepackage[spanish]{babel}
\usepackage{bbm}
\usepackage{amsthm}
\usepackage{authblk}
\usepackage{hyperref}
\usepackage[utf8]{inputenc}
\usepackage{mathrsfs}
\usepackage{amsfonts}
\usepackage{amssymb}
\usepackage{enumerate}
\usepackage[left=3cm,right=2cm,top=2cm,bottom=2cm]{geometry}
\usepackage{epsfig}
 \renewcommand{\baselinestretch}{0.93}
 \thispagestyle{empty}
 \pagestyle{empty}
\setlength{\textheight}{53pc} \setlength{\textwidth} {36pc}
\setlength{\topmargin}{-4mm}
\usepackage{multicol}
\newcommand*{\email}[1]{%
    \normalsize\href{mailto:#1}{#1}\par
    }
\setlength{\columnsep}{1cm}
\newcommand{\N}{\mathbb{N}}
\newcommand{\calM}{\cal{M}}
\newcommand{\calP}{\cal{P}}
\newcommand{\F}{\mathbb{F}}
\newcommand{\R}{\mathbb{R}}
\newcommand{\Z}{\mathbb{Z}}
\newcommand{\Q}{\mathbb{Q}}
\newcommand{\C}{\mathbb{C}}
\newcommand{\bbH}{\mathbb{H}}


\theoremstyle{definition}
\newtheorem{ejer}{Ejercicio}
\author{Marcos Morales}
\affil{Universidad Finis Terrae}
\title{Semana 4\\}



	


\begin{document}


\begin{picture}(400, 75)
   \put(-40,15){\includegraphics[width=5cm]{UFT.png}}


\end{picture}


\begin{center}
\huge{Ejercicios Control Recuperativo y Examen}
\end{center}

\begin{center}
\Large{ Marcos Morales, Camila Muñoz}
\end{center}

\begin{center}
	\large{Ecuaciones Diferenciales Ordinarias (EDO)}
\end{center}
\begin{center}
\setlength{\parindent}{0cm}\large{Clases centradas para capacitar a los estudiantes de ingeniería en Ecuaciones diferenciales ordinarias. \\}
\end{center}


\vspace{1cm}

\begin{center}
	\large{Contenidos:}
\end{center}

\setlength{\parindent}{2.95cm}- Clase de ejercicios





\vspace{6cm}


\begin{center}
\today
\end{center}
\begin{center}
Santiago, Chile
\end{center}


\pagestyle{empty}


\setcounter{page}{1}
\pagestyle{plain}
\setlength{\parindent}{0cm}





\newpage
\section{Clase de ejercicios }


\textbf{Ejercicio 1: } Encuentra una solución explícita de la siguiente ecuación diferencial de orden 1

\begin{align*}
    x\frac{dy}{dx}+y=x^2y^3
\end{align*}

Luego resuelva el PVI cuando $y(1)=-1$. \\

\textbf{Solución:} Es una ecuación de Bernoullí, usamos el cambio de variable $z=y^{1-3}=y^{-2}$, luego $z'=-2y^{-3}y'$ o bien $y'= -y^3/2 z'$, reemplazando obtenemos

\begin{align*}
    -\frac{xy^3}{2}z'+y=x^2y^3
\end{align*}

Multiplicando todo por $-2y^{-3}x^{-1}$, obtenemos

\begin{align*}
    z'-\frac{2}{x}y^{-2}= -2x
\end{align*}

o bien

\begin{align*}
    z'-\frac{2}{x}z= -2x
\end{align*}

Es una lineal, usando la fórmula general


\begin{align*}
    z &= e^{-\int \frac{-2}{x}dx} \left( \int  -2x e^{-\int \frac{2}{x}dx} dx +c \right) \\
    &= e^{\int \frac{2}{x}dx} \left( \int  -2x e^{-\int \frac{2}{x}dx} dx +c \right) \\
    &= e^{\ln(x^2)} \left( \int  -2xe^{\ln(x^{-2})} dx +c \right) \\
    &= x^2 \left( \int  -2x x^{-2}dx +c \right) \\
    &= x^2 \left( -2\int \frac{1}{x} dx +c \right) \\
    &= x^2 \left( -2\ln(x)+c \right) \\
\end{align*}

Luego

\begin{align*}
 y^{-2}&= x^2 \left( -2\ln(x)+c \right) \\
\end{align*}

Luego

\begin{align*}
 y &= \pm \frac{1}{\sqrt{x^2c-2x^2\ln(x) }}\\
\end{align*}

Ahora cambiamos $y(1)=-1$, usamos la raíz negativa, nos queda

\begin{align*}
 -1 &= \frac{-1}{\sqrt{c-2\ln(1) }}\\
\end{align*}

Luego

\begin{align*}
 \sqrt{c } &= 1\\
  c &= 1
\end{align*}

Por lo tanto la solución del PVI es

\begin{align*}
 y &= \pm \frac{1}{\sqrt{x^2-2x^2\ln(x) }}\\
\end{align*}

\newpage

\textbf{Ejercicio 2: } Resuelva la siguiente ecuación diferencial \\

\begin{align*}
    (5x+4y)dx+(4x-8y^3)dy=0
\end{align*}

\textbf{Solución: } Tenemos $M= 5x+4$ y  $N= 4x-8y^3$, luego

\begin{align*}
    M_y= 4 =N_x
\end{align*}

Por lo tanto es exacta, tenemos entonces que

\begin{align*}
    f(x,y) &= \int 5x+4y dx \\
    &= \frac{5}{2}x^2+4xy+c(y) \\
\end{align*}

Derivando respecto a $y$

\begin{align*}
    f_y &= 4x+c'(y) \\
\end{align*}

Igualando con $N$

\begin{align*}
     4x+c'(y) &= 4x-8y^3\\
     c'(y)=-8y^3
\end{align*}

Luego

\begin{align*}
    c(y ) &= \int -8y^3 dy \\
    &= -2y^4\\
\end{align*}

Tenemos entonces que

\begin{align*}
    f(x,y) &= \frac{5}{2}x^2+4xy-2y^4 \\
\end{align*}

Por lo tanto la solución de la EDO es

\begin{align*}
    \frac{5}{2}x^2+4xy-2y^4 =c \\
\end{align*}


\newpage


\textbf{Ejercicio 3: } Resuelva el siguiente PVI usando coeficientes indeterminados


\begin{align*}
    y''+y'-2y &= 1+x-x^2 \\
    y(0) &= 1 \\
    y'(0) &= 0
\end{align*}

\textbf{Solución: } Primero resolvemos la homogenea usando $y=e^{mx}$, obtenemos el polinomio asociado



\begin{align*}
    m^2+m-1 &= 0
\end{align*}

De donde se deduce que $m_1=1$ y $m_2=-2$, de esta forma la solución de la homogenea es

\begin{align*}
    y_h= c_1e^x+c_2e^{-2x}
\end{align*}

La solución particular tiene la forma 


\begin{align*}
    y_p= Ax^2+Bx+C
\end{align*}

Luego $y_p'=2Ax+B$ y $y_p''=2A$, reemplazando en la original obtenemos


\begin{align*}
    2A+2Ax+B-2(Ax^2+Bx+C) &= 1+x-x^2 \\
    -2Ax^2+2(A-B)x+2A-2C+B &= 1+x-x^2 \\
\end{align*}

De donde se deduce que $A=1/2$, luego $2(A-B)=1$ y por lo tanto $B=0$, por otra parte $2A-2C+B =1$ de donde se deduce nuevamente $C=0$, le polinomio entonces es $y_p=\frac{x^2}{2}$, de esta forma, la solución de la EDO es

\begin{align*}
    y= c_1e^x+c_2e^{-2x} +\frac{x^2}{2}
\end{align*}

Luego

\begin{align*}
    y'= c_1e^x-2c_2e^{-2x} +2x
\end{align*}

Ahora bien $y(0)=1$, obtenemos

\begin{align*}
    1= c_1+c_2
\end{align*}

Por otra parte con $y'(0)=0$, obtenemos

\begin{align*}
    0= c_1-2c_2
\end{align*}

Restando ambas ecuaciones, obtenemos

\begin{align*}
    3c_2&=1 \\
    c_2 &=\frac{1}{3}
\end{align*}

Luego

\begin{align*}
    c_1 &=\frac{2}{3}
\end{align*}

La solución del PVI es

\begin{align*}
    y= \frac{2}{3}e^x+\frac{1}{3}e^{-2x} +\frac{x^2}{2}
\end{align*}

\newpage

\textbf{Ejercicio 4: } Resuelva el siguiente PVI usando variación de parámetros


\begin{align*}
    y''+9y=\sin(3x)
\end{align*}

\textbf{Solución: } Primero resolvemos la homogenea usando $y=e^{mx}$, obtenemos el polinomio asociado



\begin{align*}
    m^2+1 &= 0
\end{align*}

De donde se deduce que $m_1=3i$ y $m_2=-3i$, $y_1=\cos(3x)$ e $y_2=\sen(3x)$ de esta forma la solución de la homogenea es

\begin{align*}
    y_h= c_1\cos(3x)+c_2\sin(3x)
\end{align*}

Ahora bien, $f(x)=\sin(3x)$, por otra parte


\begin{align*}
    W(y_1,y_2) &= \begin{vmatrix}
\cos(3x) & \sen(3x) \\
-3\sen(x) & 3\cos(x) 
\end{vmatrix} \\
&= 3\cos^2(x)+3\sen^2(x) \\
&= 3
\end{align*}

luego

\begin{align*}
    C_1(x) &= -\frac{1}{3}\int \sen(3x)\sen(3x) dx\\
    &= -\frac{1}{3}\int \sen^2(3x) dx\\
\end{align*}

haciendo $u=3x$, entonces $du=3dx$, luego $dx=du/3$

\begin{align*}
    C_1(x) &= -\frac{1}{3}\int \sen(3x)\sen(3x) dx\\
    &= -\frac{1}{9}\int \sen^2(u) du\\
    &= -\frac{1}{9}\left( \frac{u}{2}-\frac{\sen(2u)}{4} \right) \\
    &= -\frac{1}{9}\left( \frac{3x}{2}-\frac{\sen(6x)}{4} \right) \\
\end{align*}

Por otra parte

\begin{align*}
    C_2(x) &= \frac{1}{3}\int \cos(3x)\sen(3x) dx\\
\end{align*}

haciendo $u=\cos(3x)$, se tiene $du=3\sen(3x)dx$, luego

\begin{align*}
    C_2(x) &= \frac{1}{9}\int u du\\
    &= -\frac{u^2}{18}\\
    &= -\frac{1}{18}\cos^2(3x)\\
\end{align*}

Luego

\begin{align*}
    y &=   c_1\cos(3x)+c_2\sin(3x)+ \left( -\frac{1}{9}\left( \frac{3x}{2}-\frac{\sen(6x)}{4} \right) \right)\cos(3x)+ \left( -\frac{1}{18}\cos^2(3x) \right)\sen(3x) \\
    &=   c_1\cos(3x)+c_2\sin(3x)-\frac{1}{6}x\cos(3x)+\frac{1}{36}\sen(6x)\cos(3x)-\frac{1}{18}\cos^2(3x)\sen(3x) \\
    &=   c_1\cos(3x)+c_2\sin(3x)-\frac{1}{6}x\cos(3x)+\frac{1}{36}2\sen(3x)\cos(3x)\cos(3x)-\frac{1}{18}\cos^2(3x)\sen(3x) \\
     &=   c_1\cos(3x)+c_2\sin(3x)-\frac{1}{6}x\cos(3x)+\frac{1}{18}2\sen(3x)\cos^2(3x)-\frac{1}{18}\cos^2(3x)\sen(3x) \\
     &= c_1\cos(3x)+c_2\sin(3x)-\frac{1}{6}x\cos(3x)
\end{align*}




\newpage

\textbf{Ejercicio 5: } Resuelva la siguiente ecuación diferencial


\begin{align*}
    y''-5y'+6y &=\cos(t) \\
    y(0) &= 0\\
    y'(0) &= 1
\end{align*}

\textbf{Solución: }  Aplicamos Laplace

\begin{align*}
    \mathcal{L}[y'']-5\mathcal{L}[y']+6\mathcal{L}[y] &=\mathcal{L}[\cos(t)] \\
    s^2\mathcal{L}[y]-sy(0)-y'(0)-5(s\mathcal{L}[y]-y(0))+6\mathcal{L}[y] &= \frac{s}{s^2+1} \\
    \mathcal{L}[y](s^2-5s+6)-1 &= \frac{s}{s^2+1} \\
    \mathcal{L}[y](s^2-5s+6) &= \frac{s}{s^2+1}+1 \\
    \mathcal{L}[y](s^2-5s+6) &= \frac{s^2+s+1}{s^2+1} \\
    \mathcal{L}[y] &= \frac{s^2+s+1}{(s^2-5s+6)(s^2+1)} \\
    \mathcal{L}[y] &= \frac{s^2+s+1}{(s-3)(s-2)(s^2+1)} \\
\end{align*}    
    
Aplicamos fracciones parciales

\begin{align*}
    \frac{s^2+s+1}{(s-3)(s-2)(s^2+1)} &= \frac{As+B}{s^2+1}+\frac{C}{s-2}+\frac{D}{s-3} \\
    s^2+s+1 &= (As+B)(s-2)(s-3)+C(s^2+1)(s-3)+D(s^2+1)(s-2)\\
    s^2+s+1 &= (As+B)(s^2-5s+6)+C(s^3-3s^2+s-3)+D(s^3-2s^2+s-2)\\
    s^2+s+1 &= (As^3-5As^2+6As+Bs^2-5Bs+6B)+C(s^3-3s^2+s-3)+D(s^3-2s^2+s-2)\\
    s^2+s+1 &= s^3(A+C+D)+s^2(B-5A-3C-2D)+s(6A+C-5B+D)+6B-3C-2D\\
\end{align*}

Concluimos que

\begin{align*}
    A+C+D &=0 \\
    B-5A-3C-2D &=1 \\
    6A+C-5B+D&= 1 \\
    6B-3C-2D &= 1
\end{align*}

De la primera ecuación obtenemos que

\begin{align*}
    A=-C-D
\end{align*}

Reemplazando en las otras

\begin{align*}
    2C+3D+B &= 1 \\
    -5C-5D-5B &= 1 \\
    6B-3C-2D &= 1
\end{align*}

De la primera ecuación obtenemos que

\begin{align*}
    B=1-2C-3D
\end{align*}

Reemplazando en las otras dos

\begin{align*}
    -5C-5D-5(1-2C-3D) &= 1 \\
    6(1-2C-3D)-3C-2D &= 1
\end{align*}

O bien

\begin{align*}
    5C+10D&= 6 \\
    -15C-20D&= -5
\end{align*}

Dividiendo la segunda por 3

\begin{align*}
    5C+10D&= 6 \\
    -5C-\frac{20}{3}D&= -\frac{5}{3}
\end{align*}


Sumando las ecuaciones


\begin{align*}
    \frac{10}{3}D &=\frac{13}{3} \\
    D &= \frac{13}{10}
\end{align*}

Ahora reemplazando en la primera obtenemos


\begin{align*}
    5C+13&= 6 \\
    C &= -\frac{7}{5}
\end{align*}

Luego

\begin{align*}
    B &= 1+2\frac{7}{5}-3\frac{13}{10} \\
    &= 1+\frac{14}{5}-\frac{39}{10} \\
    &= -\frac{1}{10}
\end{align*}

\begin{align*}
    A &=\frac{7}{5}-\frac{13}{10} \\
    A &=\frac{1}{10}
\end{align*}

De esta forma

\begin{align*}
    \frac{s^2+s+1}{(s-3)(s-2)(s^2+1)} = \frac{s-1}{10(s^2+1)}-\frac{7}{5(s-2)}+\frac{13}{10(s-3)}
\end{align*}

Por lo tanto

\begin{align*}
    y &= \mathcal{L}^{-1}\left[ \frac{s-1}{10(s^2+1)}-\frac{7}{5(s-2)}+\frac{13}{10(s-3)}\right] (t) \\
    &= \frac{\cos(t)}{10}-\frac{\sin(t)}{10}-\frac{7}{5}e^{2t}+\frac{13}{10}e^{3t}
\end{align*}



\end{document} 
